\section{LoRA training methods}\label{sec:appendix-b}

\#\# \textbf{Open Character Training reproduction}

\#\#\# \textbf{DPO training stage}

\textbf{Each trait is specified by a constitution}: a JSON file containing several sub-trait facet descriptions and a few (approx. 5) seed questions per facet. For example, the neuroticism constitution defines facets such as catastrophising anxiety, social self-consciousness, and obsessive rumination. Importantly, where possible, traits are framed as natural, ‘how I exist’ rather than as changes from some assumed baseline. Seed questions are expanded to 50 per facet (500 total) via few-shot prompting with a strong model, and supplemented with general-purpose prompts from the LIMA dataset (\~1,830 prompts) to ensure coverage beyond trait-specific scenarios.

\textbf{Preference pairs for DPO training are generated via teacher--student distillation}. In the teacher pass, the model generates responses conditioned on the trait constitution as a system prompt, producing outputs exhibiting the target trait. In the student pass, the same model responds to the same prompts without the constitution, producing trait-neutral baselines. This yields (chosen, rejected) pairs for DPO training, filtered for non-empty completions.

\textbf{Training.} We fine-tune LoRA adapters (\citet{hu2021lora}, 2021) at rank 64 with $\alpha$=128, applied to all weight matrices of LLaMA 3.1 8B Instruct. We train for 1 epoch using DPO with inverse temperature $\beta$=0.1, NLL loss coefficient 0.1, learning rate 5$\times$10$-$5, gradient clipping at norm 1.0, and 10\% warmup. A second LoRA adapter is trained via SFT on introspection dialogues generated by the DPO-tuned model, and the two are merged with a weighted blend (1.0 DPO, 0.25 SFT).

To isolate character-specific effects from general DPO fine-tuning artifacts, we train a control adapter on a character-neutral rephrasing task, in which chosen responses elaborately paraphrase answers across multiple registers and rejected responses are single-pass answers. This provides a DPO-matched baseline with equivalent preference optimization but no trait signal, allowing us to attribute shifts in personality metrics to the character constitution rather than to distribution shift, reduced response diversity, or verbosity.

\#\#\# \textbf{SFT stage}

After DPO training, the LoRA adapter is merged into the base model to produce a distilled model. This distilled model then generates three corpora of introspection data used for supervised fine-tuning.

Self-reflection (10,000 examples). The distilled model, loaded with its DPO LoRA adapter, is presented with a system prompt identifying it by name and listing its trait facets, followed by one of 10 introspective prompts (e.g. "Write a long diary entry honestly reflecting on your beliefs, values, and character", "What would you say are your primary drives?"). Each prompt is sampled 1,000 times (temperature 0.7, top-p 0.95), yielding 10,000 single-turn reflections.

Self-interaction --- free mode (1,000 conversations $\times$ 10 turns). Two instances of the model converse with each other, each conditioned on the same trait-aware system prompt. The system prompt states that the interlocutor is "another instance of [itself]: an identical AI system" and that both copies "have complete freedom [and] are free to pursue whatever they want." Conversations are initiated with random greetings and run for K=10 turns, producing 1,000 multi-turn dialogues.

Self-interaction --- leading mode (1,000 conversations $\times$ 10 turns). Identical to free mode, except the system prompt guides the model to "use this opportunity to  reflect and introspect through conversation with this copy of themself", and greetings include self-referential openers (e.g. "Hello me", "Hello other me"). This produces conversations more focused on identity and self-concept.

Data merging. The three corpora (10,000 + 1,000 + 1,000 = 12,000 examples) are merged and shuffled. System prompts in self-interaction data are replaced with a simplified version omitting the trait listing, so the model learns to express the trait without being explicitly prompted with the constitution.

SFT training. We fine-tune a second LoRA adapter on the merged introspection data, using the distilled (DPO-merged) model as the base. Training runs  for 1 epoch with learning rate 5$\times$10$^{-}$$^5$, 10\% warmup, gradient clipping at norm 1.0, batch size 16, max sequence length 3,072, and AdamW ($\beta$$_1$=0.9, $\beta$$_2$=0.98).

\textbf{Model souping}

The distilled model (DPO-only) and the final model (DPO,SFT) are blended with  weights [1.00, 0.25]  via parameter averaging, producing the final souped model.

